%!TEX encoding = UTF-8 Unicode

\documentclass[t,compress,hyperref={hidelinks}]{beamer}

%include preamble.fmt

\newcommand{\lectureno}{IV}

\begin{document}

%include document.fmt

\title{Type theory and logic}
\subtitle{Lecture~\lectureno: meta-theoretical reasoning}
\date{4 July 2014}
\author{{\fontspec[ExternalLocation=/Library/Fonts/]{DFHeiStd-W5}\href{http://www.cs.ox.ac.uk/people/hsiang-shang.ko/}{柯向上}}}
\institute{\href{http://www.cs.ox.ac.uk/}{Department of Computer Science}\\ \href{http://www.ox.ac.uk/}{University of Oxford}\\ \href{mailto:Hsiang-Shang.Ko@@cs.ox.ac.uk}{Hsiang-Shang.Ko\char64 cs.ox.ac.uk}}

{\usebackgroundtemplate{\includegraphics[width=\paperwidth]{banner.pdf}}
\begin{frame}
\titlepage
\end{frame}}

\begin{frame}{Recap}

We have talked about the \emph{statics} and the \emph{dynamics} of type theory:
\begin{itemize}
\item the statics is typechecking, explaining why a program is a valid inhabitant of a type;
\item the dynamics is judgemental equality, capturing the notion of computation abstractly.
\end{itemize}

Both are expressed as derivations of judgements.

\end{frame}

\begin{frame}{Meta-language vs object language}

Types and programs form a language, which are talked about by a separate language of judgements and derivations.
In this case we call the former  the \emph{object language}, and the latter the \emph{meta-language}.

What we write down as types and programs are nothing more than certain syntax trees by themselves; then, at a higher level, we organise and relate these syntax trees with judgements and derivations.

Judgements and derivations can also be regarded as syntax trees to be reasoned about.
For example, consistency is a statement in which judgements and derivations are the object language and English is the meta-language.

We can also use Agda as the meta-language!

\end{frame}

\begin{frame}{Implicational fragment of propositional logic}

Today we consider only propositions formed with implication.

Each of these propositions is a finite tree whose internal nodes are implications and whose leaves are atomic propositions, which are elements of a given set |V = { A, B, C, {-"\text\ldots"-} }|.

\marker{Example} The proposition |(A -> B -> C) -> (A -> C) -> B -> C| is represented as
\[
\tikz[every node/.style={draw}, every child/.style={level distance=5.5ex}, level 1/.style={sibling distance=10em}, level 2/.style={sibling distance=5em}, level 3/.style={sibling distance=2.5em}]
  \node [circle] {$\to$}
    child { node [circle] {$\to$}
              child { node [circle] {\kern-.75pt$A$} }
              child { node [circle] {$\to$}
                        child { node [circle] {\kern-.5pt$B$} }
                        child { node [circle] {\kern-2pt$C$} } } }
    child { node [circle] {$\to$}
              child { node [circle] {$\to$}
                        child { node [circle] {\kern-.75pt$A$} }
                        child { node [circle] {\kern-2pt$C$} } }
              child { node [circle] {$\to$}
                        child { node [circle] {\kern-.5pt$B$} }
                        child { node [circle] {\kern-2pt$C$} } } };
\]

\end{frame}

\begin{frame}{Natural deduction}

Natural deduction is the type part of simple type theory (and we are considering only the implicational fragment).

\[ \AXC{$|p| \in |Gamma|$}
\rule{assum}\UIC{|Gamma ENTAILS p|} \DP \]

\[ \AXC{|Gamma , p ENTAILS q|}
\rule{$\to$I}\UIC{|Gamma ENTAILS p -> q|} \DP \]

\[ \AXC{|Gamma ENTAILS p -> q|} \AXC{|Gamma ENTAILS p|}
\rule{$\to$E}\BIC{|Gamma ENTAILS q|} \DP \]

\end{frame}

\begin{frame}{Untyped $\lambda$-calculus}


|(\x BIND x x) (\x BIND x x)|

\end{frame}

\begin{frame}{Simply typed $\lambda$-calculus (à~la Curry)}

\end{frame}

\begin{frame}{Curry--Howard isomorphism}

Derivations in natural deduction and well-typed $\lambda$-terms are in one-to-one correspondence.

That is, we can write two functions,
\begin{itemize}
\item one mapping a logical derivation in natural deduction to a $\lambda$-term and its typing derivation, and
\item the other mapping a $\lambda$-term with a typing derivation to a logical derivation in natural deduction,
\end{itemize}
and can prove that the two functions are inverse to each other.

This result is historically significant: two formalisms are developed separately from logical and computational perspectives, yet they coincide perfectly.

\end{frame}

\begin{frame}{Simply typed $\lambda$-calculus à~la Church}

The Curry--Howard isomorphism points out that derivations in natural deduction are actually $\lambda$-terms in disguise.

These $\lambda$-terms are intrinsically typed, so every term we are able to write down is necessarily well-behaved, whereas in simply typed $\lambda$-calculus à~la Curry, we can write arbitrary $\lambda$-terms, and only rule out ill-behaved ones via typing later.

\end{frame}

\begin{frame}{Semantics}

After defining a language (like the implicational fragment of propositional logic), which consists of a bunch of syntax trees, we need to specify what these trees mean.

Judgements and derivations (which form a \emph{deduction system}) assign meaning to the propositional language by specifying how it is used in formal reasoning.

We can also translate the syntax trees into entities in a well understood semantic domain.
In the case of propositional logic, we can translate propositional trees to functions on truth values. (This is the classical treatment.)

\end{frame}

\begin{frame}{Two-valued semantics of propositional logic}

\begin{itemize}

\item Define |BOOL DEFEQ { false, true }|.

\item An \emph{assignment} is a function of type |V -> BOOL|.

\item A proposition~|p| is translated into a function |(sem p) : (V -> BOOL) -> BOOL| mapping assignments to truth values.

\item An assignment~|sigma| \emph{models} a proposition~|p| exactly when |(sem p) sigma| is true, and \emph{models} a context~|Gamma| exactly when it models every proposition in~|Gamma|.

\end{itemize}

\end{frame}

\begin{frame}{Two-valued semantics of propositional logic}

\begin{itemize}

\item A proposition~|p| is \emph{satisfiable} exactly when there exists an assignment that models~|p|, and is \emph{valid} exactly when every assignment models~|p|.

\item A proposition~|p| is a \emph{semantic consequence} of a context~|Gamma| (written $|Gamma| \models |p|$) exactly when every assignment that models~|Gamma| also models~|p|.

\end{itemize}

\marker{Exercise} Show that |(p -> p -> q) -> (p -> q)| is valid for any propositions |p|~and~|q|.

\marker{Exercise} Show that a proposition~|p| is valid if and only if |p|~is a semantic consequence of the empty context.

\end{frame}

\begin{frame}{Deduction systems and semantics}

Natural deduction is \emph{sound} with respect to the two-valued semantics: whenever we can deduce $|Gamma| \vdash |p|$, it must be the case that $|Gamma| \models |p|$.

If the system of natural deduction is suitably extended, then it is also possible to prove that it is \emph{(semantically) complete} with respect to the two-valued semantics: if $|Gamma| \models |p|$, then we can construct a derivation of $|Gamma| \vdash |p|$.

\end{frame}

\end{document}
